\documentclass[11pt,a4paper]{report}

% Packages
\usepackage[utf8]{inputenc}
\usepackage[T1]{fontenc}
\usepackage{lmodern}
\usepackage[margin=2.5cm]{geometry}
\usepackage{graphicx}
\usepackage{booktabs}
\usepackage{longtable}
\usepackage{array}
\usepackage{multirow}
\usepackage{hyperref}
\usepackage{xcolor}
\usepackage{listings}
\usepackage{fancyhdr}
\usepackage{enumitem}
\usepackage{amsmath}
\usepackage{caption}
\usepackage{titlesec}

% Hyperref setup
\hypersetup{
    colorlinks=true,
    linkcolor=blue!70!black,
    urlcolor=blue!70!black,
    citecolor=green!50!black
}

% Code listing style
\lstset{
    basicstyle=\ttfamily\small,
    breaklines=true,
    frame=single,
    backgroundcolor=\color{gray!10},
    keywordstyle=\color{blue!70!black},
    commentstyle=\color{green!50!black},
    stringstyle=\color{red!70!black},
}

% Header/footer
\pagestyle{fancy}
\fancyhf{}
\fancyhead[L]{\leftmark}
\fancyhead[R]{Platelet WCM Feasibility Study}
\fancyfoot[C]{\thepage}

% Title formatting
\titleformat{\chapter}[display]
{\normalfont\huge\bfseries}{\chaptertitlename\ \thechapter}{20pt}{\Huge}
\titlespacing*{\chapter}{0pt}{-20pt}{40pt}

% Document info
\title{%
    \vspace{-2cm}
    \textbf{A Feasibility Study for Building a Whole-Cell Model of the Human Platelet}\\[1cm]
    \large Using the wcEcoli Framework as a Foundation\\[2cm]
    \includegraphics[width=0.3\textwidth]{example-image-a}\\[1cm]
}
\author{Steve Haigh}
\date{April 2026}

\begin{document}

% Title page
\begin{titlepage}
\centering
\vspace*{2cm}
{\Huge\bfseries A Feasibility Study for Building a Whole-Cell Model of the Human Platelet\par}
\vspace{1.5cm}
{\Large Using the wcEcoli Framework as a Foundation\par}
\vspace{3cm}
{\large Steve Haigh\par}
\vspace{1cm}
{\large April 2026\par}
\vfill
{\small This document synthesises the rationale, technical analysis, and implementation plan for adapting the wcEcoli whole-cell modelling framework to simulate human platelet biology.\par}
\end{titlepage}

\tableofcontents
\newpage

%==============================================================================
\part{Background: Whole-Cell Modelling Approaches}
%==============================================================================

\chapter{Introduction to Whole-Cell Modelling}

\section{What is Whole-Cell Modelling?}

Whole-cell modelling represents a paradigm shift in computational biology. Rather than modelling individual pathways or processes in isolation, a whole-cell model (WCM) integrates \emph{all} major cellular processes into a single, self-consistent simulation. The goal is to predict cellular phenotype directly from molecular-level mechanisms---to simulate the entire life cycle of a cell from first principles.

The key characteristics of whole-cell models are:

\begin{itemize}
    \item \textbf{Comprehensive scope:} All major biological processes (metabolism, gene expression, replication, etc.) are included.
    \item \textbf{Mechanistic detail:} Processes are modelled with sufficient biochemical detail to capture the underlying biology.
    \item \textbf{Integration:} Processes share a common molecular inventory and influence each other dynamically.
    \item \textbf{Emergent behaviour:} System-level properties (growth rate, cell cycle timing) emerge from the simulation rather than being imposed.
\end{itemize}

This document examines two state-of-the-art whole-cell models---wcEcoli and MC4D (the 4D minimal cell model)---and presents a feasibility study for adapting the wcEcoli framework to model human platelets.

\section{Why Model Platelets?}

Platelets are clinically critical cells involved in haemostasis, thrombosis, and inflammation. Despite extensive biochemical characterisation, no integrative whole-cell model of the platelet exists. Current computational models focus on individual pathways:

\begin{itemize}
    \item Calcium signalling models (Sveshnikova, Diamond groups)
    \item Coagulation cascade models (cascade kinetics)
    \item Integrin activation models (inside-out signalling)
\end{itemize}

A whole-cell model would enable questions that span pathways:

\begin{enumerate}
    \item How does receptor activation propagate through signalling cascades to granule release?
    \item How does platelet energetics (ATP supply) constrain activation kinetics?
    \item What explains heterogeneity in activation thresholds across individual platelets?
\end{enumerate}

These are precisely the cross-pathway, multi-scale questions that whole-cell modelling is designed to address.

%==============================================================================
\chapter{wcEcoli: The E. coli Whole-Cell Model}

\section{Overview}

wcEcoli is an open-source whole-cell computational model of \emph{Escherichia coli}, developed by the Covert Lab at Stanford University. It represents the most complete whole-cell model of a free-living organism currently available.

\subsection{Key References}

\begin{itemize}
    \item Karr JR et al. (2012) ``A whole-cell computational model predicts phenotype from genotype.'' \emph{Cell} 150(2):389--401.
    \item Macklin DN et al. (2020) ``Simultaneous cross-evaluation of heterogeneous E. coli datasets via mechanistic simulation.'' \emph{Science} 369:eaav3751.
    \item Ahn-Horst TA et al. (2022) ``An expanded whole-cell model of E. coli links cellular physiology with mechanisms of growth rate control.'' \emph{npj Syst Biol Appl} 8:30.
\end{itemize}

\subsection{What It Simulates}

The model simulates a single \emph{E. coli} cell from birth to division (60--120 minutes), tracking:

\begin{itemize}
    \item \textbf{Every gene} ($\sim$4,600 genes) and transcription state
    \item \textbf{Every mRNA molecule} ($\sim$2,000--8,000 copies) and degradation
    \item \textbf{Every protein} ($\sim$2.4 million molecules) synthesis and turnover
    \item \textbf{Every metabolite} ($\sim$1,500 species) concentrations and fluxes
    \item \textbf{DNA replication} fork positions and chromosome copy number
    \item \textbf{Cell growth} from $\sim$650 fg to $\sim$1,300 fg dry mass at division
\end{itemize}

\section{Inputs: Data and Parameter Fitting}

\subsection{Raw Data (72+ TSV Files)}

The model is parameterised from curated experimental data stored in \texttt{reconstruction/ecoli/flat/}:

\begin{table}[h]
\centering
\begin{tabular}{lll}
\toprule
\textbf{Category} & \textbf{Example Files} & \textbf{Contents} \\
\midrule
Genome & \texttt{genes.tsv}, \texttt{sequence.fasta} & 4,583 genes, coordinates \\
Transcription & \texttt{rnas.tsv}, \texttt{transcription\_factors.tsv} & RNA properties, TF binding \\
Translation & \texttt{proteins.tsv}, \texttt{translation\_efficiency.tsv} & Protein sequences, rates \\
Metabolism & \texttt{metabolic\_reactions.tsv}, \texttt{metabolism\_kinetics.tsv} & $\sim$2,500 reactions \\
Regulation & \texttt{fold\_changes.tsv}, \texttt{ppgpp\_regulation.tsv} & TF effects, stress response \\
\bottomrule
\end{tabular}
\caption{Categories of raw data used to parameterise wcEcoli.}
\end{table}

\subsection{Parameter Calculator (ParCa)}

Raw data is transformed by \texttt{fit\_sim\_data\_1.py} through a multi-stage fitting pipeline:

\begin{enumerate}
    \item \texttt{initialize()} --- Load and validate raw data
    \item \texttt{basal\_specs()} --- Calculate basal expression levels
    \item \texttt{tf\_condition\_specs()} --- Fit TF regulation parameters
    \item \texttt{fit\_condition()} --- CVXPY optimisation of expression
    \item \texttt{promoter\_binding()} --- Calculate TF--promoter binding probabilities
    \item \texttt{final\_adjustments()} --- Calibrate to target growth rate
\end{enumerate}

The output is a $\sim$180 MB serialised \texttt{sim\_data} object containing all fitted parameters.

\section{The Simulation Loop}

\subsection{Timestep Structure}

Each simulation timestep (default 1 second, dynamically adjusted 0.1--1.0s):

\begin{enumerate}
    \item \textbf{Calculate Requests:} Each process declares what molecules it needs.
    \item \textbf{Partition State:} Molecules allocated to processes by priority.
    \item \textbf{Evolve State:} Processes execute their biological function.
    \item \textbf{Merge State:} Updates combined back into global state.
    \item \textbf{Update Listeners:} Record data to disk for analysis.
\end{enumerate}

\subsection{Process Groups and Ordering}

Processes run in dependency-ordered groups:

\begin{lstlisting}[language=Python]
_processClasses = (
    (TfUnbinding,),                    # Group 1: Must run first
    (Equilibrium, TwoComponentSystem, RnaMaturation),  # Group 2
    (TfBinding,),                      # Group 3: Before transcription
    (TranscriptInitiation, PolypeptideInitiation, 
     ChromosomeReplication, ProteinDegradation,
     RnaDegradation, Complexation),    # Group 4: Initiation + decay
    (TranscriptElongation, PolypeptideElongation),  # Group 5
    (ChromosomeStructure,),            # Group 6
    (Metabolism,),                     # Group 7: Balances all changes
    (CellDivision,),                   # Group 8: Checks completion
)
\end{lstlisting}

\subsection{Resource Partitioning}

Processes compete for molecules through priority-based allocation:

\begin{itemize}
    \item \texttt{REQUEST\_PRIORITY\_METABOLISM = -10} (highest priority)
    \item \texttt{REQUEST\_PRIORITY\_DEFAULT = 0}
    \item \texttt{REQUEST\_PRIORITY\_DEGRADATION = 10} (lowest priority)
\end{itemize}

When requests exceed availability, proportional allocation with stochastic rounding is used.

\section{Major Processes}

\subsection{Transcription}

\textbf{TranscriptInitiation:} Calculates synthesis probability for each gene as:
\[
P(\text{synthesis}) = P_{\text{basal}} \times \text{TF modulation} \times \text{ppGpp effect}
\]
A multinomial draw determines which genes are initiated each timestep.

\textbf{TranscriptElongation:} Advances RNAP positions along DNA, consuming NTPs. When RNAP reaches gene end, a new RNA molecule is created.

\subsection{Translation}

Uses tRNA charging dynamics to model ribosome elongation:

\begin{enumerate}
    \item Calculate amino acid demand from sequences being translated
    \item Solve tRNA charging kinetics (steady-state approximation)
    \item Elongation rate depends on charged tRNA availability
    \item Stochastic polymerisation with resource constraints
\end{enumerate}

Key biological behaviours captured include ribosome stalling, codon-dependent rates, and competition between mRNAs.

\subsection{Metabolism}

Uses Flux Balance Analysis (FBA):

\begin{enumerate}
    \item Build stoichiometric matrix $S$ (metabolites $\times$ reactions)
    \item Set flux bounds from enzyme counts and nutrient availability
    \item Objective: maximise biomass production
    \item Solve linear program to find optimal flux distribution
\end{enumerate}

The steady-state assumption ($S \cdot v = 0$) is implicit in FBA.

\subsection{DNA Replication}

Tracks replication fork positions, with rates determined by dNTP availability. Initiation triggered by DnaA-ATP concentration threshold.

\section{Outputs}

\subsection{Directory Structure}

\begin{lstlisting}
out/{sim_dir}/
  kb/simData.cPickle           # Fitted parameters (180 MB)
  wildtype_000000/             # Variant
    000000/                    # Seed
      generation_000000/
        000000/                # Daughter
          simOut/              # Listener data
          plotOut/             # Analysis figures
\end{lstlisting}

\subsection{Listener Data (18 Listeners)}

Each listener records specific data each timestep:

\begin{table}[h]
\centering
\begin{tabular}{ll}
\toprule
\textbf{Listener} & \textbf{Data Recorded} \\
\midrule
Mass & Total mass, dry mass, mass fractions \\
BulkMolecules & Counts of all $\sim$3,000 molecule species \\
FBAResults & Reaction fluxes, exchange rates \\
GrowthLimits & Which resources limit growth \\
RibosomeData & Active ribosomes, stalling events \\
ReplicationData & Fork positions, chromosome count \\
\bottomrule
\end{tabular}
\caption{Key listeners in wcEcoli.}
\end{table}

%==============================================================================
\chapter{MC4D: The 4D Minimal Cell Model}

\section{Overview}

MC4D is a whole-cell model of JCVI-syn3A, a synthetically constructed minimal bacterium with approximately 473 genes. Published in \emph{Cell} (March 2026), it is the first whole-cell model to integrate full 3D spatial resolution with a complete biological description.

\subsection{The Organism}

JCVI-syn3A was engineered by iteratively removing genes from \emph{Mycoplasma mycoides}, retaining only those essential for growth and division. With $\sim$473 genes, it is dramatically simpler than \emph{E. coli} ($\sim$4,300 genes), making complete gene coverage tractable. However, roughly one third of genes have unknown function.

\section{The Four Computational Methods}

The model is named ``4D'' both for its spatial dimensions and for the four mathematical formalisms it integrates:

\subsection{RDME --- Reaction-Diffusion Master Equation}

The cell interior is discretised onto a 3D lattice. Molecules diffuse stochastically between voxels and react based on local concentrations. Implemented via Lattice Microbes (GPU-accelerated).

\textbf{Used for:} Translation (spatially resolved), molecular diffusion, ribosome dynamics with excluded-volume constraints.

\subsection{CME --- Chemical Master Equation}

A well-mixed stochastic formalism for globally-treated reactions where spatial homogeneity is reasonable but discrete-molecule stochasticity matters.

\textbf{Used for:} Transcription, tRNA charging.

\subsection{ODE --- Ordinary Differential Equations}

Deterministic kinetic integration of the metabolic network. Enzyme concentrations from RDME are passed to the ODE solver.

\textbf{Used for:} Metabolism, dNTP pool dynamics.

\subsection{BD --- Brownian Dynamics}

Explicit physical simulation of chromosomal DNA as a polymer, using LAMMPS molecular dynamics. SMC proteins and topoisomerases control compaction and segregation.

\textbf{Used for:} Chromosome structure and segregation.

\section{Integration Architecture}

The four methods communicate through a central Python module (\texttt{Hook.py}) that periodically:

\begin{enumerate}
    \item Extracts molecular counts from the RDME lattice
    \item Passes enzyme concentrations to ODE solver; runs metabolism
    \item Passes transcription events to CME solver
    \item Synchronises chromosome state with BD simulation
    \item Writes updated counts back to RDME lattice
\end{enumerate}

\section{Validation}

The model is validated against:

\begin{itemize}
    \item Doubling time (quantitative agreement)
    \item Origin-to-terminus ratio (DNA sequencing data)
    \item mRNA half-lives (experimental measurements)
    \item Protein spatial distributions (fluorescence data)
    \item Ribosome counts (quantitative agreement)
    \item Daughter cell heterogeneity (stochastic partitioning)
\end{itemize}

\section{Technical Requirements}

\begin{table}[h]
\centering
\begin{tabular}{lll}
\toprule
\textbf{Component} & \textbf{Software} & \textbf{Notes} \\
\midrule
RDME solver & Lattice Microbes & GPU required \\
ODE integrator & odecell & Python-based \\
Chromosome dynamics & btree\_chromo + LAMMPS & Kokkos-enabled build \\
Orchestration & Hook.py & Central integration \\
\bottomrule
\end{tabular}
\caption{Software dependencies for MC4D.}
\end{table}

%==============================================================================
\chapter{Comparison: wcEcoli vs MC4D}

\section{Side-by-Side Overview}

\begin{table}[h]
\centering
\small
\begin{tabular}{lll}
\toprule
\textbf{Dimension} & \textbf{wcEcoli} & \textbf{MC4D} \\
\midrule
Organism & \emph{E. coli} ($\sim$4,300 genes) & JCVI-syn3A ($\sim$473 genes) \\
Gene coverage & $\sim$43\% of characterised genes & $\sim$100\% (including unknowns) \\
Spatial resolution & None (well-mixed) & Full 3D (RDME lattice + BD) \\
Cell cycle & 25--130 min & $\sim$100 min \\
Metabolism & Flux Balance Analysis & ODE kinetic network \\
Transcription & Stochastic, well-mixed & CME, global stochastic \\
Translation & Stochastic, well-mixed & RDME, spatially resolved \\
Chromosome & Position tracking only & Brownian dynamics polymer \\
GPU required & No & Yes \\
Framework reuse & Explicitly designed for it & Research prototype \\
\bottomrule
\end{tabular}
\caption{Comparison of wcEcoli and MC4D.}
\end{table}

\section{Key Architectural Differences}

\subsection{Spatial Resolution}

wcEcoli treats the cell as well-mixed compartments (cytoplasm, periplasm, extracellular). MC4D places every molecule on a 3D lattice with explicit diffusion. This reveals spatial heterogeneity but requires GPU acceleration.

\subsection{Metabolic Treatment}

wcEcoli uses FBA (steady-state assumption, no kinetic parameters required). MC4D uses ODE kinetics (captures transients, requires rate constants for every reaction).

\subsection{Stochasticity}

Both models are stochastic, but MC4D more thoroughly so. All molecular processes in MC4D are governed by stochastic formalisms, producing natural population heterogeneity.

\subsection{Framework Design}

wcEcoli explicitly separates the generic engine (\texttt{wholecell/}) from organism-specific biology (\texttt{models/ecoli/}). MC4D is a research prototype not architected for reuse.

\section{Implications for Platelet Modelling}

\begin{table}[h]
\centering
\begin{tabular}{lll}
\toprule
\textbf{Consideration} & \textbf{wcEcoli} & \textbf{MC4D} \\
\midrule
Spatial heterogeneity & Limited (well-mixed) & Excellent (RDME) \\
Transient dynamics & Limited (FBA steady-state) & Excellent (ODE kinetics) \\
Chromosome modelling & Not needed for platelets & Not needed for platelets \\
Framework reusability & High & Low \\
Accessibility & Runs on laptop & Requires GPU + LAMMPS \\
\bottomrule
\end{tabular}
\caption{Relevance of each model to platelet modelling.}
\end{table}

\textbf{Pragmatic conclusion:} Build the initial platelet model using wcEcoli's framework (fast, reusable, well-tooled), using ODE kinetics for calcium/signalling rather than FBA. If spatial effects prove critical, the RDME approach from MC4D could inform a second-generation model.

%==============================================================================
\part{Technical Analysis: Framework Reusability}
%==============================================================================

\chapter{wcEcoli Architecture for Reuse}

\section{High-Level Structure}

The wcEcoli framework has four layers:

\begin{enumerate}
    \item \textbf{Reconstruction} --- builds \texttt{sim\_data} from raw data
    \item \textbf{Simulation engine} --- generic time-stepping loop
    \item \textbf{Organism model} --- processes, states, listeners (E. coli-specific)
    \item \textbf{Analysis} --- post-simulation plotting
\end{enumerate}

\section{Simulation Engine (Fully Generic)}

The base \texttt{Simulation} class (\texttt{wholecell/sim/simulation.py}) requires subclasses to define:

\begin{lstlisting}[language=Python]
class Simulation:
    _definedBySubclass = (
        "_internalStateClasses",      # State containers
        "_externalStateClasses",      # Environment
        "_processClasses",            # Biological processes
        "_initialConditionsFunction", # Set starting state
    )
    # Optional:
    _listenerClasses = ()
    _divideCellFunction = None  # Set to None for non-dividing cells
\end{lstlisting}

\textbf{Reusability: FULLY GENERIC.} No organism-specific logic.

\section{State Classes}

\subsection{BulkMolecules}

Tracks integer copy numbers in a flat array. Interface with \texttt{sim\_data}:

\begin{lstlisting}[language=Python]
sim_data.internal_state.bulk_molecules.bulk_data['id']    # molecule IDs
sim_data.internal_state.bulk_molecules.bulk_data['mass']  # masses
sim_data.molecule_groups.bulk_molecules_binomial_division # division lists
\end{lstlisting}

\textbf{Reusability: HIGH.} One hardcoded reference to ATP[c] needs generalisation.

\subsection{UniqueMolecules}

Tracks individual molecules with per-instance attributes (structured numpy arrays).

\textbf{Reusability: MEDIUM.} Division mode loading is E. coli-specific; subclass for platelets (or use empty division modes since platelets don't divide).

\subsection{LocalEnvironment}

Tracks external concentrations with media timeline support. Maps well to platelet agonist addition experiments.

\textbf{Reusability: MEDIUM-HIGH.}

\section{Process Architecture}

The base \texttt{Process} class (\texttt{wholecell/processes/process.py}) is fully generic:

\begin{lstlisting}[language=Python]
class Process:
    def initialize(self, sim, sim_data): pass  # Setup
    def calculateRequest(self): pass           # Declare needs
    def evolveState(self): pass                # Execute biology
    def isTimeStepShortEnough(self, dt, safety): return True
\end{lstlisting}

E. coli processes are \textbf{NOT reusable}---they serve only as examples.

\section{Reusability Summary}

\begin{table}[h]
\centering
\small
\begin{tabular}{llll}
\toprule
\textbf{Component} & \textbf{Generic?} & \textbf{Platelet Action} \\
\midrule
Simulation engine & YES & Use as-is, subclass \\
Process base class & YES & Use as-is \\
View system & YES & Use as-is \\
BulkMolecules & MOSTLY & Override ATP[c] hardcode \\
UniqueMolecules & MOSTLY & Subclass, custom division modes \\
LocalEnvironment & MOSTLY & Simplify exchange logic \\
Containers & YES & Use as-is \\
Listener base & YES & Use as-is \\
Cell division & NO & Not needed (platelets don't divide) \\
E. coli processes & NO & Write new platelet processes \\
E. coli reconstruction & NO & Write new platelet reconstruction \\
\bottomrule
\end{tabular}
\caption{Component reusability for platelet model.}
\end{table}

\section{Minimum sim\_data Interface}

For a minimal platelet simulation, \texttt{SimulationDataPlatelet} must provide:

\begin{lstlisting}[language=Python]
class SimulationDataPlatelet:
    # Required by InternalState:
    submass_name_to_index = {'protein': 0, 'lipid': 1, 'small_molecule': 2}
    compartment_abbrev_to_index = {'c': 0, 'dg': 1, 'ag': 2, 'e': 3}
    
    # Required by BulkMolecules:
    constants.n_avogadro
    internal_state.bulk_molecules.bulk_data['id']
    internal_state.bulk_molecules.bulk_data['mass']
    molecule_groups.bulk_molecules_binomial_division  # empty list OK
    
    # Required by UniqueMolecules:
    internal_state.unique_molecule.unique_molecule_definitions  # can be {}
    internal_state.unique_molecule.unique_molecule_masses
    
    # Required by LocalEnvironment:
    external_state.make_media
    external_state.saved_media
\end{lstlisting}

%==============================================================================
\part{The Platelet Model}
%==============================================================================

\chapter{Biological Context}

\section{Key Differences from E. coli}

\begin{table}[h]
\centering
\begin{tabular}{lll}
\toprule
\textbf{Feature} & \textbf{E. coli} & \textbf{Platelet} \\
\midrule
Nucleus & Yes & No (anucleate) \\
Transcription & Active & None (minimal residual) \\
Replication & Continuous & None \\
Division & Binary fission & None \\
Lifespan & Indefinite (growth) & $\sim$7--10 days \\
Primary function & Growth \& reproduction & Haemostasis \& signalling \\
Time axis & Birth $\to$ division (60 min) & Resting $\to$ activated (seconds--minutes) \\
\bottomrule
\end{tabular}
\caption{Key biological differences between E. coli and platelets.}
\end{table}

\section{Platelet-Specific Biology to Model}

\subsection{Compartmentalisation}

Platelets have highly organised internal compartments:

\begin{itemize}
    \item \textbf{Dense granules ($\sim$3--8 per platelet):} Contain ADP, ATP, Ca$^{2+}$, serotonin
    \item \textbf{Alpha granules ($\sim$40--80 per platelet):} Contain fibrinogen, vWF, growth factors
    \item \textbf{Open canalicular system:} Membrane invaginations for secretion
    \item \textbf{Mitochondria:} ATP production, calcium buffering
    \item \textbf{Cytosol:} Signalling cascades, cytoskeleton
\end{itemize}

\subsection{Receptors and Signalling}

Major platelet receptors:

\begin{itemize}
    \item \textbf{GPVI:} Collagen receptor (ITAM signalling)
    \item \textbf{PAR1/PAR4:} Thrombin receptors (GPCR)
    \item \textbf{P2Y1/P2Y12:} ADP receptors (GPCR)
    \item \textbf{$\alpha$IIb$\beta$3:} Fibrinogen receptor (integrin)
\end{itemize}

All converge on calcium mobilisation as the central integrator.

\subsection{Calcium Dynamics}

Ca$^{2+}$ signalling involves:

\begin{enumerate}
    \item Receptor activation $\to$ PLC$\gamma$2 $\to$ IP$_3$ production
    \item IP$_3$ $\to$ ER calcium release (STIM1/Orai1)
    \item Store-operated calcium entry (SOCE)
    \item Mitochondrial Ca$^{2+}$ uptake (MCU) and release (NCLX)
    \item Plasma membrane Ca$^{2+}$ ATPase (PMCA) efflux
\end{enumerate}

\subsection{Granule Exocytosis}

Calcium-dependent secretion:

\begin{itemize}
    \item SNARE machinery (VAMP-3/8, syntaxin-11, SNAP-23)
    \item Different Ca$^{2+}$ thresholds for $\alpha$- vs $\delta$-granules
    \item Positive feedback via released ADP
\end{itemize}

\subsection{Metabolism}

Platelets rely on:

\begin{itemize}
    \item Glycolysis (primary ATP source at rest)
    \item Oxidative phosphorylation (mitochondria)
    \item Metabolic switching upon activation (PDK/PDH)
\end{itemize}

%==============================================================================
\chapter{Literature Analysis}

\section{Key Data Sources}

Based on analysis of $\sim$90 papers across four Zotero collections:

\subsection{Calcium Dynamics}

\textbf{High Priority:}

\begin{itemize}
    \item \textbf{Dolan \& Diamond (2014):} 34-species ODE model of Ca$^{2+}$ signalling---the most comprehensive implementable model
    \item \textbf{Kleppe (2018):} ODE model of NO/cGMP/cAMP inhibitory pathway (missing from other Ca$^{2+}$ models)
    \item \textbf{Ghatge (2026):} MCU (mitochondrial calcium uniporter) quantification
    \item \textbf{Fernández (2021):} Ca$^{2+}$ threshold for TMEM16F/anoctamin-6 activation
\end{itemize}

\subsection{Receptor Signalling}

\textbf{High Priority:}

\begin{itemize}
    \item \textbf{Dunster (2015):} Data-driven ODE model of GPVI $\to$ Syk $\to$ PLC$\gamma$2 pathway
    \item \textbf{Mazet (2020):} ODE model of PI cycle (PIP$_2$ $\to$ IP$_3$ + DAG)
\end{itemize}

\subsection{Granule Exocytosis}

\textbf{High Priority:}

\begin{itemize}
    \item \textbf{Fitch-Tewfik \& Flaumenhaft (2013):} SNARE machinery and fusion pore dynamics
    \item \textbf{Lopez et al. (2015):} Ca$^{2+}$ thresholds for $\alpha$- vs $\delta$-granule secretion
\end{itemize}

\subsection{Metabolism}

\textbf{High Priority:}

\begin{itemize}
    \item \textbf{Thomas et al. (2014):} iAT-PLT-636 constraint-based metabolic reconstruction (636 reactions)
    \item \textbf{Flora et al. (2023):} PDK/PDH metabolic switch mechanism
\end{itemize}

\subsection{Proteomics}

\textbf{Critical Resource:}

\begin{itemize}
    \item \textbf{Burkhart et al. (2012):} Quantitative platelet proteome---copy numbers for $\sim$4,000 proteins
\end{itemize}

%==============================================================================
\chapter{Implementation Plan}

\section{Phased Approach}

\subsection{Phase 1: Framework Validation (2--3 weeks)}

\textbf{Goal:} Confirm wcEcoli engine runs a minimal non-E. coli simulation.

\begin{enumerate}
    \item Clone and run standard E. coli simulation end-to-end
    \item Map the architecture (processes, states, data flow)
    \item Document minimum \texttt{sim\_data} interface requirements
\end{enumerate}

\textbf{Deliverable:} Architecture documentation (complete---see Part II).

\subsection{Phase 2: Minimal Platelet Stub (3--4 weeks)}

\textbf{Goal:} Create a runnable ``blank cell'' in the wcEcoli framework.

\begin{enumerate}
    \item Create \texttt{models/platelet/} namespace with:
    \begin{itemize}
        \item \texttt{sim/simulation.py} --- \texttt{PlateletSimulation} class
        \item \texttt{sim/initial\_conditions.py} --- Set starting state
        \item One toy process (e.g., \texttt{RestingDecay})
        \item One listener (e.g., \texttt{Mass})
    \end{itemize}
    \item Create \texttt{reconstruction/platelet/} with:
    \begin{itemize}
        \item \texttt{simulation\_data.py} --- \texttt{SimulationDataPlatelet}
        \item Minimal molecule definitions
        \item Compartment definitions (cytosol, granules, extracellular)
    \end{itemize}
    \item Verify simulation runs without E. coli dependencies
\end{enumerate}

\textbf{Deliverable:} Runnable platelet stub ($\sim$500 lines new code).

\subsection{Phase 3: Biological Processes (8--12 weeks)}

\subsubsection{Phase 3a: Calcium Dynamics (3--4 weeks)}

Implement ODE-based calcium process:

\begin{itemize}
    \item ER calcium stores (IP$_3$R, SERCA)
    \item Mitochondrial buffering (MCU, NCLX)
    \item Plasma membrane fluxes (PMCA, SOCE)
    \item Base on Dolan \& Diamond 34-species model + Kleppe inhibitory pathway
\end{itemize}

\textbf{Validation:} Reproduce published Ca$^{2+}$ transients for thrombin stimulation.

\subsubsection{Phase 3b: Granule Release (2--3 weeks)}

Implement granule exocytosis:

\begin{itemize}
    \item Granules as unique molecules with state (loaded/fusing/released)
    \item Ca$^{2+}$-dependent fusion rate function
    \item Cargo release into extracellular compartment
    \item Different thresholds for $\alpha$- vs $\delta$-granules
\end{itemize}

\textbf{Validation:} Reproduce secretion kinetics from Lopez et al. (2015).

\subsubsection{Phase 3c: Receptor Signalling (2--3 weeks)}

Implement receptor $\to$ Ca$^{2+}$ cascade:

\begin{itemize}
    \item One receptor pathway initially (GPVI or thrombin/PAR)
    \item Connect to IP$_3$ production
    \item Trigger calcium dynamics from receptor occupancy
\end{itemize}

\textbf{Validation:} Dose-response curves for agonist $\to$ Ca$^{2+}$ mobilisation.

\subsubsection{Phase 3d: Metabolism (2--3 weeks)}

Implement simplified metabolic process:

\begin{itemize}
    \item ATP production (glycolysis + OXPHOS)
    \item ATP consumption by processes
    \item Metabolic switch upon activation
\end{itemize}

Consider adapting iAT-PLT-636 reconstruction (Thomas et al. 2014).

\subsection{Phase 4: Integration and Validation (4+ weeks)}

\begin{enumerate}
    \item Cross-validate processes (Ca$^{2+}$ drives granule release drives ADP feedback)
    \item Compare to experimental time courses
    \item Sensitivity analysis on key parameters
    \item Document limitations and future directions
\end{enumerate}

\section{v0.1 Target Definition}

A concrete, falsifiable deliverable:

\begin{quote}
\textbf{Simulate thrombin-induced Ca$^{2+}$ rise and dense granule release over 5 minutes. Plot secretion kinetics and compare to published data.}
\end{quote}

This requires:
\begin{itemize}
    \item Working simulation engine with platelet namespace
    \item Calcium dynamics process (ODE)
    \item Granule release process
    \item Thrombin receptor signalling (minimal)
    \item Listeners for Ca$^{2+}$ trace and granule counts
\end{itemize}

\section{Proposed File Structure}

\begin{lstlisting}
models/platelet/
    sim/
        simulation.py           # PlateletSimulation(Simulation)
        initial_conditions.py   # Set initial state
    processes/
        calcium_dynamics.py     # Ca2+ store release, buffering
        granule_release.py      # Exocytosis
        receptor_signalling.py  # GPVI, PAR, P2Y12 -> Ca2+
        metabolism.py           # ATP production/consumption
    listeners/
        mass.py
        activation_state.py
        granule_counts.py
        calcium_trace.py

reconstruction/platelet/
    simulation_data.py          # SimulationDataPlatelet
    knowledge_base_raw.py       # Literature values
    dataclasses/
        constants.py
        molecule_groups.py
        state/
            internal_state.py
            external_state.py
\end{lstlisting}

%==============================================================================
\chapter{Feasibility Assessment}

\section{Technical Feasibility}

\subsection{Framework Suitability}

The wcEcoli framework is well-suited for adaptation:

\begin{itemize}
    \item \textbf{Modular design:} Clear separation between engine and organism-specific code
    \item \textbf{Documented interfaces:} State and process APIs are well-defined
    \item \textbf{Proven at scale:} Handles 3,000+ molecules, 15+ processes
    \item \textbf{Existing tooling:} Analysis scripts, data I/O, logging infrastructure
\end{itemize}

\subsection{Biological Compatibility}

Key adaptations required:

\begin{itemize}
    \item \textbf{No division:} Set \texttt{\_divideCellFunction = None} (straightforward)
    \item \textbf{ODE processes:} Replace FBA with kinetic ODEs for signalling (requires new process implementations)
    \item \textbf{Time axis:} Seconds--minutes activation rather than hours growth (adjust \texttt{lengthSec})
\end{itemize}

\subsection{Risk Assessment}

\begin{table}[h]
\centering
\begin{tabular}{llll}
\toprule
\textbf{Risk} & \textbf{Likelihood} & \textbf{Impact} & \textbf{Mitigation} \\
\midrule
Framework too E. coli-specific & Low & High & Phase 2 stub validates genericity \\
Insufficient parameter data & Medium & Medium & Use Burkhart proteomics + literature \\
ODE stiffness issues & Medium & Low & Use adaptive solvers (scipy) \\
Scope creep & High & Medium & Strict v0.1 target definition \\
\bottomrule
\end{tabular}
\caption{Risk assessment for platelet model development.}
\end{table}

\section{Scientific Feasibility}

\subsection{Data Availability}

Platelet biology is well characterised:

\begin{itemize}
    \item \textbf{Proteomics:} Burkhart et al. provides copy numbers for $\sim$4,000 proteins
    \item \textbf{Signalling kinetics:} Multiple ODE models available (Dolan, Dunster, Kleppe)
    \item \textbf{Metabolomics:} Metabolic reconstructions exist (Thomas et al.)
    \item \textbf{Validation data:} Extensive experimental literature on activation kinetics
\end{itemize}

\subsection{Model Scope}

A proof-of-concept model does not require:

\begin{itemize}
    \item Complete coverage of all pathways (start with 1--2 receptors)
    \item Spatial resolution (well-mixed is adequate initially)
    \item Full metabolic network (simplified ATP module sufficient)
\end{itemize}

\section{Resource Requirements}

\subsection{Computational}

\begin{itemize}
    \item \textbf{Development:} Standard laptop sufficient
    \item \textbf{Simulation:} Expected $<$10 minutes per 5-minute simulation
    \item \textbf{No GPU required} (unlike MC4D)
\end{itemize}

\subsection{Time Estimate}

\begin{table}[h]
\centering
\begin{tabular}{lll}
\toprule
\textbf{Phase} & \textbf{Duration} & \textbf{Cumulative} \\
\midrule
Phase 1: Framework validation & 2--3 weeks & 3 weeks \\
Phase 2: Minimal stub & 3--4 weeks & 7 weeks \\
Phase 3a: Calcium dynamics & 3--4 weeks & 11 weeks \\
Phase 3b: Granule release & 2--3 weeks & 14 weeks \\
Phase 3c: Receptor signalling & 2--3 weeks & 17 weeks \\
Phase 3d: Metabolism & 2--3 weeks & 20 weeks \\
Phase 4: Integration & 4+ weeks & 24 weeks \\
\bottomrule
\end{tabular}
\caption{Estimated timeline for proof-of-concept model.}
\end{table}

\section{Conclusion}

Building a whole-cell model of the human platelet using the wcEcoli framework is \textbf{technically feasible} and \textbf{scientifically justified}:

\begin{enumerate}
    \item The wcEcoli framework provides a mature, well-tested foundation with explicit support for reuse.
    \item Platelet biology is well characterised with sufficient data for parameterisation.
    \item A proof-of-concept model (Ca$^{2+}$ dynamics + granule release) is achievable in $\sim$6 months.
    \item The model would fill a gap in the field---no integrative WCM of platelets currently exists.
\end{enumerate}

The key success factors are:

\begin{itemize}
    \item \textbf{Ruthless scope control:} Start with minimal biology, validate, then extend
    \item \textbf{ODE-first approach:} Use kinetic models for signalling rather than FBA
    \item \textbf{Clear validation targets:} Define falsifiable predictions from the outset
\end{itemize}

%==============================================================================
\part{Appendices}
%==============================================================================

\appendix

\chapter{Key Literature References}

\section{Whole-Cell Modelling}

\begin{enumerate}
    \item Karr JR et al. (2012) A whole-cell computational model predicts phenotype from genotype. \emph{Cell} 150(2):389--401.
    \item Macklin DN et al. (2020) Simultaneous cross-evaluation of heterogeneous E. coli datasets. \emph{Science} 369:eaav3751.
    \item Thornburg ZR et al. (2026) Bringing the genetically minimal cell to life on a computer in 4D. \emph{Cell}.
\end{enumerate}

\section{Platelet Calcium Signalling}

\begin{enumerate}
    \item Dolan AT, Diamond SL (2014) Systems modeling of Ca2+ homeostasis and mobilization in platelets. \emph{PLoS ONE}.
    \item Kleppe LS (2018) Mathematical modelling of NO/cGMP/cAMP signalling in platelets.
    \item Ghatge M et al. (2026) The mitochondrial calcium uniporter regulates calcium dynamics. 
\end{enumerate}

\section{Platelet Receptor Signalling}

\begin{enumerate}
    \item Dunster JL et al. (2015) Regulation of early steps of GPVI signal transduction by phosphatases. \emph{PLoS Comput Biol}.
    \item Mazet F et al. (2020) A model of the PI cycle reveals regulating roles of lipid-binding proteins.
\end{enumerate}

\section{Platelet Metabolism}

\begin{enumerate}
    \item Thomas SG et al. (2014) Network reconstruction of platelet metabolism. \emph{Blood}.
    \item Flora GD et al. (2023) Mitochondrial PDK2/4 contribute to platelet function.
\end{enumerate}

\section{Platelet Proteomics}

\begin{enumerate}
    \item Burkhart JM et al. (2012) The first comprehensive and quantitative analysis of human platelet protein composition. \emph{Blood} 120(15):e73--e82.
\end{enumerate}

\chapter{Glossary}

\begin{description}
    \item[FBA] Flux Balance Analysis --- constraint-based metabolic modelling
    \item[RDME] Reaction-Diffusion Master Equation --- spatially-resolved stochastic simulation
    \item[CME] Chemical Master Equation --- well-mixed stochastic simulation
    \item[ODE] Ordinary Differential Equations --- deterministic kinetic modelling
    \item[ParCa] Parameter Calculator --- wcEcoli's data fitting pipeline
    \item[GPVI] Glycoprotein VI --- platelet collagen receptor
    \item[PAR] Protease-Activated Receptor --- thrombin receptor
    \item[SOCE] Store-Operated Calcium Entry
    \item[MCU] Mitochondrial Calcium Uniporter
    \item[SNARE] Soluble NSF Attachment Protein Receptor --- membrane fusion machinery
\end{description}

\end{document}
